\renewcommand{\tCoverLine}{The arcade game on MSX}
\renewcommand{\tComputerText}{\textbf{MSX1 or later}, with a slot for the game cartridge and another for the V9990.}
\renewcommand{\tMemoryText}{\textbf{64 KB of RAM recommended}.}
\renewcommand{\tVideoText}{A \textbf{Yamaha V9990} cartridge is required. Start-up credits use the MSX's own video output. The cover, title and game use the V9990 output.}
\renewcommand{\tSoundText}{The MSX's built-in PSG. No extra sound cartridge is needed.}
\renewcommand{\tControlsText}{One player. Use the keyboard or a two-button joystick in port 1.}
\renewcommand{\tSetGame}{Insert the \textbf{\gameTitle{} cartridge} in another free slot.}
\renewcommand{\tSetOn}{Turn on the MSX. Its own video output shows the credits and \texttt{V9990 found!}. If it says \texttt{V9990 NOT found!}, switch off, check the card and try again.}
\renewcommand{\tBlackScreen}{\textbf{If the game screen stays black}, check the MSX video output for the start-up message, then switch the monitor to the V9990 output. In openMSX, attach the V9990 extension and load the ROM as ASCII16.}
\renewcommand{\tCreditsTitle}{9. Credits and rights}
\renewcommand{\tSourceNote}{Source code available on GitHub. \textbf{Please improve the game if you feel so!}}
\newcommand{\tCreditsExtra}{\item Cover and label use the project's opening illustration. The screenshots are from the MSX port.}
\newcommand{\tRelease}{Version}
